\section*{Appendix P. Strong-start research and representation optimization}
\label{app:strong-start}
\renewcommand{\thetable}{P.\arabic{table}}
\renewcommand{\theHtable}{P.\arabic{table}}
\setcounter{table}{0}

\paragraph{Separating the starting design from the research controller.}
We first conduct a researcher-designed training/development optimization study, then give its improved starting designs identically to both research controllers. The engineering search is not credited as native harness output. Its 794 candidates comprise 540 DTI retrieval/fusion configurations, 72 proteomics feature configurations, and 49/97/36 Reverse-ID/phenocopy/CEX policies. Each task's source is selected using training/development evidence before that batch applies the selected candidate to the historical final role. Previously published tests remain exposed; these are versioned development studies, not new blind benchmarks.

\paragraph{Stronger DTI starting point.}
The selected retrieval predictor uses all 50,149 training pairs. Morgan-fingerprint and protein-composition cosine similarities receive weights 0.25 and 0.75; the 50 nearest pairs receive softmax weights with temperature 0.05. Its probability is mixed equally with TAPB. The training representations and original source-row tie order are fixed. Selection maximizes the lower AUROC across two SMILES-hash validation halves, with mean-half AUROC as tie-breaker, using the mean validation predictions of the five existing TAPB fits. Each deployment applies this single selected rule to one TAPB fit, not an ensemble of five models. Their mean random/unseen-drug/unseen-protein AUROCs are 0.9710/0.9588/0.9317. A separately labeled five-model ensemble scores 0.9769/0.9671/0.9484 and has additional inference cost; it is not substituted for the single-model comparisons.

\paragraph{Stronger row-local proteomics starting point.}
The 72 feature configurations vary six temporal representations, four drug weights and three output scales. Five compound-grouped training folds plus one full-training fit give 432 fits. Selection maximizes the mean of pooled OOF AP and development AP. The chosen map uses the observed 24h profile, with independently normalized protein/drug blocks, drug weight 0.5, and overall scale $3/\sqrt{1.25}$; the logistic objective remains $C=1$. Scaling changes effective regularization. This deterministic predictor attains final AP/AUROC/BCE 0.3699/0.7302/0.3741. Feature discovery is centralized engineering; deployment is separately verified in a fresh coordinator process forbidden from reading biological arrays. Ten local workers reproduce centralized predictions within $1.23\times10^{-10}$.

\paragraph{Fixed strong-start native-controller matrix.}
Twenty campaigns compare direct generation and the pinned production loop, five pairs per task (742--746), with identical task-specific starts, numerical APIs, selectors and aggregate feedback. Each has ceilings of three candidate opportunities, four evaluations including baseline, and 24 LLM calls. Actual calls and unused opportunities are reported; equal ceilings do not imply equal inference spend. DTI research uses the fixed TAPB342 and weighted-retrieval validation predictions; selected programs deploy to TAPB342--346 by the fixed trajectory-minus-400 mapping. Both visible-half AUROCs must improve by more than $10^{-4}$. PTPC uses the original three training-compound folds and ten-client partition; both pooled OOF AP and minimum-fold AP must improve by more than $10^{-4}$. The 12 charged fold-fit ceiling, local logistic objective, and gradient RPC remain unchanged. No final metric enters either controller's prompt or selector. All twenty campaign selections close before their final deployment evaluation.

The task adapter changes agenda/entry routing without changing production-driver instruction bytecode; synthetic tests check route binding and exact numerical starting points. An auxiliary file initially shadowed Python's standard-library \texttt{queue} module in eighteen process launches. It was renamed before those launches performed scientific work or model calls; their original logs remain preserved. This was not a score-driven restart. The two unaffected campaigns continued unchanged.

\begin{table}[ht]
\centering\small\setlength{\tabcolsep}{3pt}
\begin{tabularx}{\linewidth}{@{}Xrrrr@{}}
\toprule
Method & Random & Unseen drug & Unseen protein & PTPC AP \\
\midrule
Fixed predictor & $0.9545\,\pm\,0.0083$ & $0.9328\,\pm\,0.0148$ & $0.8868\,\pm\,0.0362$ & $0.2783$ \\
Shared strong start & $\textbf{0.9710}\,\pm\,0.0034$ & $\textbf{0.9588}\,\pm\,0.0061$ & $\textbf{0.9317}\,\pm\,0.0153$ & $0.3699$ \\
Harness-free & $\textbf{0.9710}\,\pm\,0.0034$ & $\textbf{0.9588}\,\pm\,0.0061$ & $\textbf{0.9317}\,\pm\,0.0153$ & $0.3716\,\pm\,0.0055$ \\
Our research loop & $\textbf{0.9710}\,\pm\,0.0034$ & $\textbf{0.9588}\,\pm\,0.0061$ & $\textbf{0.9317}\,\pm\,0.0153$ & $\textbf{0.3727}\,\pm\,0.0039$ \\
\bottomrule\end{tabularx}
\caption{Strong-start comparison. DTI columns are AUROC; PTPC uses primary AP. Mean $\pm$ sample SD across five fixed trajectories; deterministic PTPC controls have no SD. Fixed predictors are TAPB (DTI) and logistic regression (PTPC). Both research controllers receive the same improved starting source and the same opportunity/call ceilings; actual inference costs differ. Bold marks the best displayed mean, including display ties.}
\label{tab:strong-start}

\end{table}

All twenty scientific campaigns complete. Every DTI selection retains its shared starting combiner, making direct and loop predictions identical under each paired backbone fit. PTPC loop/direct mean AP is 0.372710/0.371641. There are 141 successful LLM calls: DTI direct/loop use 15/45 and PTPC direct/loop use 15/66. Their respective input/output-token totals are 237,835/25,617; 844,746/39,468; 243,994/27,716; and 1,291,055/62,754. Final deployment yields eleven PTPC labels from eight exact-source fits and 45 DTI labels. Saved predictions reproduce the original scorers to $10^{-12}$; ten native-loop persistence histories and all twenty source closures are verified. Complete scalar records, sample SDs, source hashes and costs are provided alongside the tables.

\paragraph{Paired primary comparisons.}
The primary contrasts are loop minus direct in unseen-drug AUROC and final PTPC AP. Each uses all five declared pairs, all 32 two-sided sign flips and Holm correction across the two task-level contrasts. These pairs describe research trajectories on fixed test data, not independent institutions; the smallest attainable unadjusted two-sided value is 0.0625.
\begin{center}\small\begin{tabular}{lrrr}\toprule
Task & Mean loop--direct & Exact $p$ & Holm $p$ \\
\midrule
DTI & +0.000000 & 1.0000 & 1.0000 \\
PTPC & +0.001069 & 0.6250 & 1.0000 \\
\bottomrule\end{tabular}\end{center}

\begingroup\footnotesize\setlength{\tabcolsep}{3pt}
\begin{longtable}{llrrrrr}
\caption{All strong-start DTI metrics; five fixed pairs per method.}\label{tab:strong-dti-all}\\
\toprule Split & Method & AUROC & AP & MCC & Acc. & BCE \\\midrule\endfirsthead
\toprule Split & Method & AUROC & AP & MCC & Acc. & BCE \\\midrule\endhead
random & Fixed predictor & 0.9545 & 0.9492 & 0.7907 & 0.8948 & 0.3557 \\
random & Shared strong start & \textbf{0.9710} & \textbf{0.9709} & \textbf{0.8255} & \textbf{0.9126} & \textbf{0.2331} \\
random & Harness-free & \textbf{0.9710} & \textbf{0.9709} & \textbf{0.8255} & \textbf{0.9126} & \textbf{0.2331} \\
random & Our research loop & \textbf{0.9710} & \textbf{0.9709} & \textbf{0.8255} & \textbf{0.9126} & \textbf{0.2331} \\
unseen\_drug & Fixed predictor & 0.9328 & 0.9416 & 0.7364 & 0.8705 & 0.4342 \\
unseen\_drug & Shared strong start & \textbf{0.9588} & \textbf{0.9670} & \textbf{0.7819} & \textbf{0.8929} & \textbf{0.2730} \\
unseen\_drug & Harness-free & \textbf{0.9588} & \textbf{0.9670} & \textbf{0.7819} & \textbf{0.8929} & \textbf{0.2730} \\
unseen\_drug & Our research loop & \textbf{0.9588} & \textbf{0.9670} & \textbf{0.7819} & \textbf{0.8929} & \textbf{0.2730} \\
unseen\_protein & Fixed predictor & 0.8868 & 0.8573 & 0.6537 & 0.8254 & 0.5885 \\
unseen\_protein & Shared strong start & \textbf{0.9317} & \textbf{0.9353} & \textbf{0.7155} & \textbf{0.8574} & \textbf{0.3625} \\
unseen\_protein & Harness-free & \textbf{0.9317} & \textbf{0.9353} & \textbf{0.7155} & \textbf{0.8574} & \textbf{0.3625} \\
unseen\_protein & Our research loop & \textbf{0.9317} & \textbf{0.9353} & \textbf{0.7155} & \textbf{0.8574} & \textbf{0.3625} \\
\bottomrule\end{longtable}
\begin{longtable}{lrrrr}
\caption{All strong-start PTPC metrics. Higher is better except BCE; bold includes all best displayed ties. AP is the primary metric.}\label{tab:strong-ptpc-all}\\
\toprule Method & AP & AUROC & BCE & Acc. \\\midrule\endfirsthead
\toprule Method & AP & AUROC & BCE & Acc. \\\midrule\endhead
Fixed predictor & $0.2783$ & $0.7273$ & $0.3905$ & $\textbf{0.8514}$ \\
Shared strong start & $0.3699$ & $\textbf{0.7302}$ & $\textbf{0.3741}$ & $0.8446$ \\
Harness-free & $0.3716\,\pm\,0.0055$ & $0.7246\,\pm\,0.0034$ & $0.3780\,\pm\,0.0023$ & $\textbf{0.8514}\,\pm\,0.0000$ \\
Our research loop & $\textbf{0.3727}\,\pm\,0.0039$ & $0.7281\,\pm\,0.0041$ & $0.3757\,\pm\,0.0028$ & $0.8473\,\pm\,0.0037$ \\
\bottomrule\end{longtable}

\endgroup

\paragraph{Fixed first-index native trajectory.}
For PTPC742, the starting pooled/minimum-fold OOF AP is 0.430994/0.383892. The three retained revisions score 0.449933/0.386053, 0.450078/0.387487 and 0.450192/0.387899. The final source independently normalizes the observed 24h profile $p$, the 24h-minus-6h difference $r$, and compound vector $d$. With $s$ the row mean of $r$, it constructs $x=[p,0.5r,0.5d(1+0.75s)]$ and returns $3x/\max(\lVert x\rVert_2,10^{-12})$. The fitted coefficients are learned through the unchanged local-gradient service. Its final AP is 0.375813; same-index direct AP is 0.375886. The case demonstrates successful source revision on visible evidence, not a selected win over direct or a causal biological mechanism.

\paragraph{Additional local-data availability.}
Holding the improved PTPC map and optimizer fixed, five predeclared compound partitions each expose the first one, three or ten clients. Their training-condition counts are 44--59, 130--159 and 491. Mean final AP is 0.3088/0.3183/0.3699, AUROC 0.6901/0.7033/0.7302 and BCE 0.4015/0.3877/0.3741. All fifteen fits use a raw-array-denied coordinator and data-owning workers. Ten-client partitions contain the identical full pool, so their repeated AP is numerical replay, not five independent laboratory experiments. The measured benefit is access to additional local data; no intrinsic optimization advantage over centralized fitting, differential privacy or secure aggregation is claimed.

\paragraph{Non-improving optimization routes are retained.}
A separate 288-configuration DTI probe replaces protein composition with frozen ESM2 residue-mean embeddings, testing training-only centering, two drug similarities, three block weights, four neighbor counts, two temperatures and three mixture weights. Its validation-selected candidate reaches AUROC 0.9669/0.9474/0.9086, below the composition-retrieval starting point. It is not adopted. A second cell study starts from the historical winning formulas and uses both original visible roles: 45 Mahalanobis variants, 48 calibrated phenocopy policies and 90 CEX acquisition policies. Selected final scores are 0.803922/0.723634/0.093995, versus unchanged starting scores 0.803922/0.723917/0.096127. No improved cell-policy claim is made. Complete arm/metric and cost records accompany the strong-start table; none of these engineering outcomes are relabeled as a controller comparison.
